\section{Формальная проверка геометрии и её применение к наблюдателю}\label{app:lean}

\subsection{Две ветви доказательства}
Файл \texttt{lean/Rule110Exact.lean} импортирует \texttt{Finite} и \texttt{Late}. В пространстве имён \texttt{Rails} корневой результат имеет имя \texttt{essential\_iff}. Его левая часть, \texttt{EssentialAt}\,$t\,j$, требует двух конфигураций, совпадающих во всех позициях, кроме, возможно, $j$, и дающих различные значения целевой клетки в момент $t$. Различие выходов вынуждает различие входов в $j$. В правой части Lean сначала вычисляет $t-t/3$ в натуральных числах с целочисленным делением, а затем вкладывает результат в $\mathbb Z$.

В прямом направлении \texttt{left\_dependency\_bound} и \texttt{right\_dependency\_bound} исключают существенность координаты за пределами окна. В обратном направлении сертификат указывает для каждого $(t,j)$ допустимую пару траекторий. Для $t<2485$ файл \texttt{lean/Finite.lean} применяет \texttt{checkSound} к конечным блокам, выделяет свидетельствующий бит в позиции физического дефекта $s=-j$ и переносит свидетельство на прямую. Генератор \texttt{scripts/generate.py} предлагает битовые снимки. Совместимость соседних блоков и покрытие проверяются выражениями Lean до вывода леммы \texttt{Rails.Finite.finite\_all}.

Для $t\ge2485$ файл \texttt{lean/Late.lean} сводит координату к одной из пяти полос подраздела~\ref{sec:five-bands}. Используются полулучевой сертификат ширины 106 из \texttt{Rule110LeftCheckpoints}, эфирный клин \texttt{EtherWedge}, доказательство \texttt{Crystal50WedgeAll}, узкая полоса \texttt{LeftCrystal50} и правая граничная полоса ширины 20 из \texttt{Rule110Right20}. Файл \texttt{Crystal50WedgeAll.lean} доказывает сохранение одновременно кристаллической среды и 16-битного окна при обоих значениях следующего внешнего бита. Конечные таблицы переходов и условие склейки проверяются в построении общего инварианта. Из покрытия всех пяти случаев следует \texttt{Rails.late\_all}.

\subsection{Область машинного свидетельства}
При подготовке монографии корень импортов заново собран на macOS arm64 с Lean~4.35.0, указанным в \texttt{lean-toolchain}. \src{verification/monograph/geometry-axioms.log}{Журнал проверки геометрии} фиксирует обработку 263 модулей и покрытие всех $5\,147\,263$ пар конечной области в 100 блоках. Для \texttt{Rails.late\_all} и \texttt{Rails.essential\_iff} ядро вывело \texttt{[propext, Classical.choice, Quot.sound]}. Те же списки зафиксированы в \src{verification/v2/final.log}{журнале предыдущего выпуска}. Конечный критерий проверки персистентной существенности на произвольном конечном кольце установлен отдельно в предложении~\ref{prop:ring_cert}. Для кольца из пяти клеток контрпример дан в замечании~\ref{rem:ring5}. Числа для кольца из 20 клеток получены полным конечным перебором.

Команда \texttt{make verify} заново генерирует конечный сертификат, проверяет замыкание импортов \texttt{Rule110Exact} в порядке зависимостей и сверяет строку об аксиомах корневой теоремы. Команда \texttt{make verify-series} включает также проверку модуля наблюдателя и реестров.

\subsection{Односторонние автономные тайлы и две бесконечные приграничные полосы}\label{app:half-ray-tiles}

Из Леммы~\ref{lem:three_step} композицией вытекает, что область зависимости $3m$-тактного выхода в ячейке $s$ удовлетворяет включению $\mathrm{Dep}_{3m}(s) \subseteq [s - 2m, s + 3m]$ для любых конфигураций.

\begin{lemma}[Одностороннее полулучевое замыкание]\label{lem:one_sided}
Пусть $A_t = F^t(A_0)$, $p > 0$ и $a(t + p) = a(t) + \delta$. Если на базовых классах $t_0 \le t < t_0 + p$ выполнено $A_{t+p}(s + \delta) = A_t(s)$ для всех $s \ge a(t)$, и вычисление $A_{t+p}$ на полулуче $[a(t+p), \infty)$ зависит только от значений $A_t$ на $[a(t), \infty)$, то сдвиговое равенство $A_{t+p}(s + \delta) = A_t(s)$ сохраняется для всех $t \ge t_0$ и $s \ge a(t)$. Аналогичное утверждение верно для левых полулучей $s \le a(t)$.
\end{lemma}

Применяя Лемму~\ref{lem:one_sided} к правому физическому фронту $e(t) = L(t)$ с шагом $(\Delta t, \Delta s) = (30, +20)$ (где $e(t+30) = e(t) + 20$, а область зависимости нового тайла ширины $W$ лежит в $[e(t) - W + 1, e(t) + 50]$ и не затрагивает неизвестную внутреннюю область $s < e(t) - W + 1$) и аналогично к левому фронту $s = -t$ с шагом $(\Delta t, \Delta s) = (40,-40)$, получаем первую крупную бесконечную теорему.

\begin{theorem}[Две приграничные полосы]\label{thm:dual_strips}
Для каждого целого $t \ge 300$ выполнено
\begin{equation}\label{eq:dual_strips}
[-L(t), -L(t) + 105] \cap \mathbb{Z} \;\cup\; [t - 7, t] \cap \mathbb{Z} \;\subseteq\; D_t.
\end{equation}
\end{theorem}

Левая полоса ширины $106$ клеток сертифицируется $20$ фазами фонов $(0001111001)$, $(1011)$ и $(001011)$ с глубиной тайла $167$ клеток на $30$ базовых тактах $t = 300, \dots, 329$. Правая полоса ширины $8$ клеток сертифицируется шестью фазами с шагом $(\Delta t, \Delta s) = (40,-40)$ на базовых тактах $t = 300, \dots, 339$. Происхождение промежуточных состояний удостоверено по схеме \emph{proof-carrying computation} на битовых векторах \texttt{BitVec 1320} с проверкой каждого перехода $F_{1320}^{h_i}(z_i) = z_{i+1}$ ($h_i \in \{1, 10\}$) в ядре Lean 4.35.0.

\subsection{Баллистические лучи и границы локальных сертификатов}

Второй механизм неограниченного продолжения во времени основан на сохранении различия вдоль дискретных мировых линий.

\begin{lemma}[Сохранение правого полулуча]\label{lem:ballistic_ray}
Зафиксируем границу полулуча $s_0\in\mathbb Z$. Если для траектории $A$ в момент $T$ выполнено $\forall s \ge s_0,\; A_{T+3}(s + 2) = A_T(s)$, то $\forall q \ge 0,\; \forall s \ge s_0,\; A_{T+3q}(s + 2q) = A_T(s)$. В частности, если это условие выполнено одновременно для $B$ и $\widetilde B$ и $B_T(s_0) \ne \widetilde B_T(s_0)$, то $B_{T+3q}(s_0 + 2q) \ne \widetilde B_{T+3q}(s_0 + 2q)$ для всех $q \ge 0$.
\end{lemma}

Для двух арифметических серий якорей на фонах $(0000101111)$ и $(001011)$ при $0 \le n \le 32$:
\begin{equation}
T_A(n) = 594 + 882n,\quad s_A(n) = 256 + 378n; \qquad
T_B(n) = 884 + 882n,\quad s_B(n) = 380 + 378n,
\end{equation}
вычислительно сертифицированы $66$ начальных точек, порождающих бесконечные баллистические лучи существенности на постоянных вдоль каждого луча глубинах $d_A(n) = 140 + 210n$ и $d_B(n) = 210 + 210n$ от фронта $e(t) = L(t)$.

В последующих модулях \texttt{StaircaseA} и \texttt{StaircaseB} формулы обеих серий доказаны для произвольных $n,q\ge0$. Для этого три локально согласованные строки склеиваются в точную эволюцию. На передних периодических средах серий A и B обнаружено $45\,255$ покрытых классов из $66\,150$; с третьей серией --- $53\,550$. Число классов относится к таблице остатков $315\times210$. Применение таблицы в физической точке требует также проверки границ соответствующего клина.

Важную методологическую роль сыграли строгие отрицательные результаты, уточнившие границы применимости простых гипотез:
\begin{itemize}
    \item \textbf{Опровержение локально-тайловой макроиндукции по $n \mapsto n+1$:} при переходе на макрошаг $(\Delta T, \Delta s) = (882, +378)$ универсальная область зависимости составляет $\mathrm{Dep}_{882}(s) \subseteq [s - 588, s + 882]$ (прообраз $[s - 210, s + 1260]$). Прямыми контрпримерами в \texttt{check\_macro\_boundary.py} показано, что фиксация локальных тайлов радиуса $64$ (для $T=594$) или $186$ (для $T=884$) недостаточна: изменение клеток вне тайла меняет целевой бит через $882$ такта.
    \item \textbf{Неполнота наивных угловых секторов:} для фона $(0001111001)$ наблюдаемый до $t = 2000$ левый сектор со скоростями $[-0.66, -0.64]$ обнаруживает первый пропуск на более длинном горизонте.
    \item \textbf{Разделение периодичности и полноты покрытия:} глубокая сдвиговая рекуррентность ($261$ клетка для $(0000001001)$) не гарантирует сплошного покрытия без объединения нескольких фаз по всем классам вычетов.
\end{itemize}

\subsection{Конструктивное замыкание бесконечного левого клина: теорема \texttt{wedge\_all}}

Чтобы преодолеть несоизмеримость скоростей между левым краем светового конуса (скорость $-1$), фоновым кристаллом периода $50$ и фронтом распространения возмущения, вводится дискретная параметризация чётных тактов $n \in \mathbb{N}$:
\begin{equation}
\mathrm{tm}(n) = 2450 + 2n, \qquad a(n) = -350 - n.
\end{equation}
Здесь правая граница клина $a(n)$ смещается влево со скоростью $-1/2$ клетки за такт ($1$ клетка за $2$ такта). Для каждой из шести фаз $p \in \{\mathrm{p0}, \mathrm{p1}, \mathrm{p2}, \mathrm{p3}, \mathrm{p7}, \mathrm{p9}\}$ пишем $\widetilde B_p(t,s)=\widetilde B^p_t(s)$. Сопряжённый инвариант на слое $n$ объединяет совпадение с фоновым кристаллом $\mathrm{Cm}_p$ на интервале слева от $a(n)$ и принадлежность $16$-битного граничного буфера узлу конечного графа переходов:
\begin{align}
\mathrm{Medium}(p,n)\ \iff\ {}&
\forall s\in[-\mathrm{tm}(n)+1300,a(n))\cap\mathbb{Z},\notag\\
&\widetilde B_p(\mathrm{tm}(n),s)=\mathrm{Cm}_p(\mathrm{tm}(n),s), \\
\mathrm{Window}(p,n,i)\ \iff\ {}&
\forall k\in\{0,\dots,15\},\notag\\
&\widetilde B_p(\mathrm{tm}(n),a(n)+k)=\mathrm{testBit}(\mathrm{word}(p,n,i),k).
\end{align}

\begin{lemma}[16-битный детерминированный трансдьюсер и склейка со швом]\label{lem:transducer}
Для каждой фазы $p$ таблица \texttt{Crystal50Packed} задаёт $50$-периодический граф $36$-битных записей (младшие $16$ бит кодируют окно $\mathrm{word}(p, n, i)$, биты $16\dots 25$ и $26\dots 35$ — переходы при правом бите $r \in \{0, 1\}$).
\begin{enumerate}[label=(\alph*)]
    \item На шаге $n$ $4$-битный левый контекст $\mathrm{ctx}(p, n \bmod 50)$ в точках $a(n) - 4 + m$ ($m \in \{0, 1, 2, 3\}$) имеет глубину $t + s = 2096 + n + m$ и принадлежит уже доказанному интервалу $\mathrm{Medium}(p, n)$.
    \item Двукратное побитовое сжатие $\mathrm{shrink}(\mathrm{shrink}(\mathrm{input}(p, n, i, r)))$ вычисляет $17$ бит конфигурации в момент $\mathrm{tm}(n+1)$ на $[a(n)-2, a(n)+14]$: младший бит ($k=0$) даёт $\mathrm{Cm}_p(\mathrm{tm}(n+1), a(n) - 2)$, а старшие $16$ бит ($k=1\dots 16$) совпадают с узлом $z$ следующего слоя графа в позиции $a(n+1) = a(n) - 1$.
    \item Левый край полосы $-\mathrm{tm}(n+1) + 1300 \le s < -\mathrm{tm}(n) + 1302$ имеет инвариантную глубину $k = s + \mathrm{tm}(n+1) \in \{1300,1301,1302,1303\}$ и закрывается теоремой \texttt{seam} (\texttt{Crystal50Seam}) благодаря инвариантности $t + s$ при сдвиге $(t, s) \mapsto (t + 50, s - 50)$.
\end{enumerate}
\end{lemma}

\begin{theorem}[\texttt{Rails.Crystal50Wedge.wedge\_all}]\label{thm:wedge_all}
Для любого натурального $t \ge 2485$ и любого целого $s$, удовлетворяющего неравенствам
\begin{equation}\label{eq:wedge_region}
-t + 1302 \le s \le 59 - \lfloor t/2 \rfloor,
\end{equation}
выполняется предикат \texttt{EssentialAt}\,$t\,(-s)$. Верхней границы по времени нет.
\end{theorem}

\begin{proof}[Схема доказательства]
Положим $n = \lfloor (t - 2450) / 2 \rfloor$. При $t \ge 2485$ имеем $n \ge 17$, а правая граница $59 - \lfloor t/2 \rfloor = -1166 - n = a(n) - 816$. Для чётных тактов $t = \mathrm{tm}(n)$ равенство $\widetilde B_p(t, s) = \mathrm{Cm}_p(t, s)$ следует непосредственно из инварианта $\mathrm{Medium}(p, n)$ для всех шести фаз $p$. Для нечётных тактов $t = \mathrm{tm}(n) + 1$ одношаговый конус зависимости $[s - 1, s + 1]$ целиком лежит внутри $[-\mathrm{tm}(n) + 1300, a(n))$, и совпадение переносится леммой локальности \texttt{carry}. Применение моста \texttt{essential\_of\_cm} завершает доказательство.
\end{proof}

\subsection{Полное покрытие на бесконечной прямой}\label{sec:five-bands}
В координате входа $j=-s$ положим $b=\lfloor t/2\rfloor$. При $t\ge2485$ последовательно применяются следующие пять результатов. Обе границы в каждой строке включены.
\begin{center}
\begin{tabular}{lll}
\toprule
Начало & Конец & Основание \\
\midrule
$-L(t)$ & $-L(t)+105$ & Полулучевой сертификат ширины 106 \\
$-L(t)+106$ & $b-60$ & \texttt{EtherWedge.half\_sector} \\
$b-59$ & $t-1302$ & \texttt{Crystal50Wedge.wedge\_all} \\
$t-1301$ & $t-20$ & \texttt{LeftCrystal50.left\_strip} \\
$t-19$ & $t$ & \texttt{right\_boundary\_width20} \\
\bottomrule
\end{tabular}
\end{center}
Три фазы эфира периода 14 дают вторую строку уже при $t\ge210$. Последняя строка усиливает раннюю правую полосу ширины восемь. При $t=2485$ третья строка состоит из единственной клетки $j=1183$. Начало каждой следующей строки на единицу больше конца предыдущей. Лемма \texttt{Rails.late\_all} объединяет пять доказательств.

Первые численные проверки охватили $0\le t\le2000$. Позднее эксперимент с фиксированными фазами семи слов охватил времена до $t=5000$. Для доказательства достаточно конечной части, примыкающей к пяти полосам.

\paragraph{Конечный сертификат.}
Для $0\le t<2485$ берём все циклические фазы семи слов. Их длины и целые кодировки (младший бит соответствует первой клетке) равны
\[
(8,244),\ (10,632),\ (10,976),\ (10,576),\ (4,13),\ (6,52),\ (6,48).
\]
Фоны представлены векторами длины 5520; однобитный дефект внесён в позицию 2760. Периоды всех семи слов делят 5520. Для каждой строки времени строим побитовое объединение различий всех пар и проверяем, что оно содержит маску физического интервала $[-t,L(t)]$. Это $\sum_{t=0}^{2484}(2t+1-\lfloor t/3\rfloor)=5\,147\,263$ пары. При $s$ внутри интервала выполнено $-5520<s-t$ и $s+t<5520$. Поэтому лемма \texttt{local\_obs} переносит результат на бесконечную прямую с единственным переворотом.

Начальная строка вычисляется из семи слов. Генератор предлагает промежуточные снимки с шагом 25. Lean проверяет маски каждого блока и равенство его конечного состояния начальному снимку следующего. Последний блок содержит десять строк. Индукция по длине блока связывает бит маски с конкретной парой траекторий. Таким образом, лемма \texttt{Rails.Finite.finite\_all} доказывает существенность для каждого входа конечной области.

\begin{proof}[Завершение доказательства теоремы~\ref{thm:gap-free}]
Включение $D_t\subseteq I_t$ следует из теоремы~\ref{thm:upper_bound}. Для обратного включения применяем конечный сертификат при $t<2485$ и пять полос при $t\ge2485$. Каждый свидетель различия в физической координате $s$ переносится сдвигом в свидетельство существенности входа $-s$.
\end{proof}

\subsection{Переход к неполному наблюдению}
В \texttt{observer/Rule110Observer.lean} импортируются \texttt{Rule110Exact} и \texttt{Rule110FrontGlue}. В доказательстве \texttt{essential\_iff\_inCone} корневая теорема применяется к смещению $j-s$, после чего сдвиг конфигурации переносит свидетельство к целевой клетке $s$. Тем самым точное окно цели равно
\[
  [s-L(t),s+t]\cap\mathbb Z.
\]
Определение \texttt{CanPredict} ставит следующий вопрос. Достаточна ли последовательность считываний с заранее заданных сайтов до установленного срока, чтобы вычислить целевой бит для любого допустимого входа? Лемма \texttt{canPredict\_iff} выражает достаточность постоянством ответа на слоях одинаковой наблюдаемой истории. В файле также заданы контрпримеры к достаточности одного геометрического покрытия. Это разграничивает достижимость информации по конусу и её фактическую доступность данному наблюдателю.

Исходник наблюдателя проверен на Lean~4.35.0 с обоими каноническими импортами. В \src{verification/monograph/observer-axioms.log}{журнале наблюдателя} присутствуют все 11 из 11 ответов на \texttt{\#print axioms}, от \texttt{canPredict\_iff} до \texttt{coverage\_does\_not\_suffice}. Для каждой теоремы ядро указало ровно \texttt{[propext, Classical.choice, Quot.sound]}. Числа OLMo относятся к отдельным опытам раздела~\ref{sec:olmo}.
